\section{งานวิจัยและทฤษฎีที่เกี่ยวข้อง}
\label{sec:review}

\subsection{ทฤษฎีที่เกี่ยวข้อง}

\subsubsection{มาตรวัดผลลัพธ์ที่รายงานโดยผู้ป่วยและแบบประเมิน DASH และ TESS}

แบบประเมิน Disabilities of the Arm, Shoulder and Hand (DASH) พัฒนาขึ้นเพื่อวัดความบกพร่องและอาการของรยางค์บนจากมุมมองของผู้ป่วย \citep{Hudak1996}
ส่วนแบบประเมิน Toronto Extremity Salvage Score (TESS) พัฒนาขึ้นเพื่อวัดความสามารถทางกายของผู้ป่วยเนื้องอกกระดูกและเนื้อเยื่ออ่อนภายหลังการผ่าตัดแบบสงวนรยางค์
โดยแยกเป็นฉบับรยางค์บนและรยางค์ล่าง \citep{Davis1996} ข้อคำถามส่วนใหญ่ของแบบประเมินทั้งสองถามถึงระดับความยากลำบากในการทำกิจกรรมในชีวิตประจำวัน
และทุกข้อมีตัวเลือกคำตอบ 5 ระดับที่มีคำบรรยายกำกับ ดังรูปที่~\ref{fig:dash-tess}
ทั้งนี้ DASH มีข้อคำถามส่วนหนึ่งที่ถามถึงความรุนแรงของอาการและผลกระทบต่อการทำงานและกิจกรรมทางสังคม \citep{Hudak1996}
ส่วน TESS มีตัวเลือก ``ไม่เกี่ยวข้อง'' (not applicable) สำหรับกิจกรรมที่ผู้ป่วยไม่ได้ทำในชีวิตประจำวัน \citep{Davis1996}

\begin{figure}[htbp]
\centering
\includegraphics[width=0.6\textwidth]{figures/dash_example.png}
\caption{ตัวอย่างคำถามและตัวเลือกคำตอบของแบบประเมิน DASH ฉบับภาษาไทย}
\label{fig:dash-tess}
\end{figure}

คะแนนรายข้อ (1--5) ของทั้งสองแบบประเมินถูกแปลงเป็นคะแนนรวมในช่วง 0--100 ตามสมการที่~\eqref{eq:dash_score} และ~\eqref{eq:tess_score} ตามลำดับ
\begin{equation}
\mathrm{DASH} = \left( \frac{1}{n} \sum_{i=1}^{n} x_i - 1 \right) \times 25
\label{eq:dash_score}
\end{equation}
\begin{equation}
\mathrm{TESS} = \frac{\sum_{i=1}^{n} x_i - n}{4n} \times 100
\label{eq:tess_score}
\end{equation}
โดยที่ $x_i$ คือคะแนนของข้อที่ $i$ และ $n$ คือจำนวนข้อที่ตอบ
DASH ต้องตอบอย่างน้อย 27 จาก 30 ข้อจึงจะคำนวณคะแนนได้ \citep{IWH2010DASHScoring} ส่วน TESS ไม่นำข้อที่ตอบว่าไม่เกี่ยวข้องมาคำนวณ \citep{CODE2025TESS}
คะแนน DASH ที่สูงหมายถึงความบกพร่องมาก ในขณะที่คะแนน TESS ที่สูงหมายถึงความสามารถทางกายที่ดี

แบบประเมินทั้งสองมีฉบับภาษาไทยที่ผ่านการตรวจสอบความตรงและความเที่ยงแล้ว \citep{Buntragulpoontawee2018, Srisawat2018}
อย่างไรก็ตาม ผลการตรวจสอบนี้ใช้ได้กับแบบประเมินที่มีถ้อยคำและรูปแบบตามที่ใช้ในการตรวจสอบเท่านั้น
การนำแบบประเมินไปใช้ในรูปแบบอิเล็กทรอนิกส์จึงต้องคงถ้อยคำคำถามและตัวเลือกคำตอบตามต้นฉบับ
เพื่อรักษาความเท่าเทียมของการวัด (measurement equivalence) \citep{Coons2009} งานวิจัยนี้จึงออกแบบระบบให้ถามข้อคำถามด้วยถ้อยคำตามต้นฉบับทุกครั้ง

\subsubsection{การรู้จำเสียงพูดอัตโนมัติ}

การรู้จำเสียงพูดอัตโนมัติ (Automatic Speech Recognition: ASR) คือการแปลงสัญญาณเสียงพูดเป็นข้อความ
แบบจำลองที่ใช้อย่างแพร่หลายในปัจจุบัน ได้แก่ Whisper \citep{Radford2023} ซึ่งเป็น Transformer แบบตัวเข้ารหัส-ตัวถอดรหัส (encoder-decoder)
ที่รับสเปกโตรแกรม log-Mel ของสัญญาณเสียงช่วงละ 30 วินาที และสร้างข้อความทีละโทเค็นแบบถดถอยในตัว (autoregressive)
งานวิจัยนี้ใช้ Typhoon Whisper Turbo \citep{TyphoonWhisperTurbo} ซึ่งปรับจูนจาก Whisper Large-v3 Turbo
ที่ลดจำนวนชั้นของตัวถอดรหัสจาก 32 ชั้นเหลือ 4 ชั้น และมีพารามิเตอร์ประมาณ 809 ล้าน \citep{OpenAI2024WhisperTurbo}
โดยฝึกด้วยข้อมูลเสียงภาษาไทยประมาณ 11,000 ชั่วโมง ซึ่งข้อความถอดเสียงผ่านการปรับให้อยู่ในรูปคำอ่านมาตรฐาน (text normalization)
เพื่อแก้ความกำกวมของการเขียนภาษาไทย เช่น ตัวเลขที่อ่านได้หลายแบบและไม้ยมก \citep{Sirichotedumrong2026}

เนื่องจากภาษาไทยไม่มีการเว้นวรรคระหว่างคำ การวัดอัตราความผิดพลาดระดับคำจึงต้องตัดคำก่อน
ASR ภาษาไทยจึงนิยมประเมินด้วยอัตราความผิดพลาดระดับอักขระ (Character Error Rate: CER) ซึ่งเท่ากับ $(S + D + I)/N$
โดยที่ $S$, $D$ และ $I$ คือจำนวนอักขระที่ถูกแทนที่ ถูกลบ และถูกแทรก ตามลำดับ และ $N$ คือจำนวนอักขระในข้อความอ้างอิง \citep{Sirichotedumrong2026}
Typhoon Whisper Turbo มีค่า CER ร้อยละ 4.79 บนชุดข้อมูลเสียงทั่วไป Gigaspeech2 และร้อยละ 6.85 บนชุดข้อมูลเสียงที่มีสัญญาณรบกวน TVSpeech
ใกล้เคียงกับ Typhoon Whisper Large-v3 ที่ฝึกด้วยข้อมูลชุดเดียวกัน (ร้อยละ 4.69 และ 6.32) \citep{Sirichotedumrong2026}
แบบจำลองจึงเหมาะกับการประมวลผลบนเซิร์ฟเวอร์ภายใน และเหมาะกับการถอดเสียงแบบออฟไลน์มากกว่าแบบต่อเนื่องตามเวลาจริง (streaming) \citep{TyphoonWhisperTurbo}
ซึ่งสอดคล้องกับระบบที่บันทึกคำตอบของผู้ป่วยทีละข้อก่อนถอดเสียง
ทั้งนี้ ผลลัพธ์ของแบบจำลองอยู่ในรูปคำอ่าน เช่น ``สิบ'' แทน ``10'' \citep{Sirichotedumrong2026} ซึ่งส่วนตีความคำตอบต้องรองรับ

\subsubsection{การสังเคราะห์เสียงพูด}

การสังเคราะห์เสียงพูด (Text-to-Speech: TTS) คือการแปลงข้อความเป็นเสียงพูด
OmniVoice \citep{Zhu2026OmniVoice} เป็นแบบจำลอง TTS แบบ zero-shot ที่รองรับมากกว่า 600 ภาษา รวมถึงภาษาไทย
โดยใช้ Transformer แบบสองทิศทาง (bidirectional) แปลงข้อความเป็นโทเค็นเสียงโดยตรงแบบไม่ถดถอยในตัว (non-autoregressive)
แบบจำลองถูกฝึกด้วยข้อมูลเสียงจากแหล่งเปิด 581,000 ชั่วโมง ซึ่งเป็นภาษาไทยประมาณ 10,500 ชั่วโมง
และเสียงภาษาไทยที่สร้างขึ้นมีค่า CER ร้อยละ 1.23 บนชุดทดสอบ FLEURS ใกล้เคียงกับเสียงจริง (ร้อยละ 1.45) \citep{Zhu2026OmniVoice}
งานวิจัยนี้ใช้ OmniVoice สร้างเสียงคำถาม เสียงขอให้ตอบใหม่ และเสียงอ่านตัวเลือกไว้ล่วงหน้า และตรวจฟังก่อนนำไปใช้
เพื่อให้ผู้ป่วยทุกคนได้ยินถ้อยคำและการออกเสียงเดียวกัน

\subsubsection{แบบจำลองภาษาขนาดใหญ่และการถอดรหัสแบบมีข้อจำกัด}
\label{sec:qwen35}

Qwen3.5 \citep{QwenTeam2026Qwen35} เป็นแบบจำลองภาษาขนาดใหญ่แบบเปิดเผยค่าน้ำหนัก (open-weight) ที่รองรับ 201 ภาษา รวมถึงภาษาไทย
โดยมีแบบหนาแน่น (dense) ขนาด 0.8--27 พันล้านพารามิเตอร์ และแบบผสมผู้เชี่ยวชาญ (Mixture-of-Experts: MoE) ขนาด 35--397 พันล้านพารามิเตอร์
และเลือกเปิดหรือปิดโหมดคิด (thinking mode) ผ่านพารามิเตอร์ของแม่แบบการสนทนาได้ \citep{QwenTeam2026Qwen35Card}
แบบจำลองแบบ open-weight สามารถติดตั้งบนเซิร์ฟเวอร์ภายในสถานพยาบาลได้ จึงไม่ต้องส่งข้อมูลของผู้ป่วยไปยังบริการภายนอก \citep{Wiest2024}
และการที่ Qwen3.5 มีหลายขนาดทำให้เลือกขนาดให้เหมาะกับหน่วยความจำของเซิร์ฟเวอร์ที่มีอยู่ได้

แบบจำลองภาษาขนาดใหญ่สร้างข้อความทีละโทเค็นโดยเลือกโทเค็นถัดไปจากการแจกแจงความน่าจะเป็นของทุกโทเค็นในคลังคำศัพท์
ผลลัพธ์จึงอาจมีรูปแบบไม่ตรงตามที่โปรแกรมต้องการ
การถอดรหัสแบบมีข้อจำกัด (constrained decoding) แก้ปัญหานี้โดยกำหนดรูปแบบผลลัพธ์ด้วยโครงสร้างข้อมูล เช่น JSON schema
แล้วในแต่ละขั้นจะตัดโทเค็นที่ทำให้ผลลัพธ์ผิดรูปแบบออกก่อนเลือกโทเค็นถัดไป
ผลลัพธ์จึงอยู่ในรูปแบบที่กำหนดเสมอ \citep{Willard2023} และวิธีนี้ช่วยให้แบบจำลองทำงานที่ต้องการผลลัพธ์แบบมีโครงสร้างได้โดยไม่ต้องปรับจูน \citep{Geng2023}

\subsubsection{ตัวชี้วัดสำหรับการประเมินผล}
\label{sec:metrics}

งานวิจัยนี้ประเมินการจับคู่คำตอบเข้ากับระดับคะแนนใน 2 ระดับ ระดับแรกคือระดับรายข้อ
ซึ่งเทียบระดับคะแนนที่ระบบเลือกกับระดับคะแนนอ้างอิง ได้แก่ ระดับคะแนนเป้าหมายของข้อมูลสังเคราะห์หรือคะแนนที่พยาบาลให้
โดยใช้ accuracy และ weighted kappa ระดับที่สองคือระดับคะแนนรวม ซึ่งเทียบคะแนนรวมที่คำนวณจากคะแนนรายข้อของระบบ
กับคะแนนรวมที่คำนวณจากคะแนนอ้างอิงด้วยสูตรเดียวกัน โดยใช้ intraclass correlation coefficient

\paragraph{ความถูกต้อง (accuracy)}
คือสัดส่วนของตัวอย่างที่ทำนายระดับคะแนนตรงกับระดับคะแนนอ้างอิง ($N_{\text{correct}}/N$) ซึ่งแปลผลได้ง่าย แต่มีข้อจำกัด 2 ประการ
ประการแรก นับการทำนายผิดทุกกรณีเท่ากัน จึงไม่สะท้อนว่าการทำนายผิดไปยังระดับที่อยู่ติดกันควรถือว่าผิดน้อยกว่าการทำนายผิดไปยังระดับที่อยู่ห่างออกไป
ประการที่สอง อาจสูงเกินจริงเมื่อจำนวนตัวอย่างของแต่ละระดับไม่สมดุล (class imbalance)

\paragraph{ค่าสัมประสิทธิ์แคปปาแบบถ่วงน้ำหนัก (weighted kappa)}
คือสถิติที่วัดความสอดคล้องระหว่างผู้ให้คะแนน 2 ฝ่ายสำหรับข้อมูลเชิงอันดับ (ordinal data) \citep{Cohen1968}
ซึ่งชดเชยข้อจำกัดทั้งสองประการของ accuracy โดยให้น้ำหนักความไม่สอดคล้องตามระยะห่างระหว่างระดับ และหักส่วนที่สอดคล้องกันโดยบังเอิญออก ดังสมการที่~\eqref{eq:kappa}
\begin{equation}
\kappa_w = 1 - \frac{\sum_{i=1}^{C}\sum_{j=1}^{C} w_{ij}\, o_{ij}}{\sum_{i=1}^{C}\sum_{j=1}^{C} w_{ij}\, e_{ij}}, \qquad w_{ij} = \frac{(i-j)^2}{(C-1)^2}
\label{eq:kappa}
\end{equation}
โดยที่ $o_{ij}$ คือสัดส่วนของตัวอย่างที่ฝ่ายที่ 1 ให้ระดับ $i$ และฝ่ายที่ 2 ให้ระดับ $j$
$e_{ij}$ คือสัดส่วนที่คาดหวังโดยบังเอิญ ซึ่งเท่ากับผลคูณของสัดส่วนที่แต่ละฝ่ายให้ระดับนั้น
$w_{ij}$ คือน้ำหนักความไม่สอดคล้องแบบกำลังสอง (quadratic weight) และ $C$ คือจำนวนระดับ
ค่า $\kappa_w$ เท่ากับ 1 เมื่อสอดคล้องกันทั้งหมด และเท่ากับ 0 เมื่อสอดคล้องกันเท่ากับที่คาดหวังโดยบังเอิญ
โดยแปลผลตามเกณฑ์ของ \citet{Landis1977} ดังตารางที่~\ref{tab:agreement-interp}

\paragraph{ค่าสัมประสิทธิ์สหสัมพันธ์ภายในชั้น (intraclass correlation coefficient: ICC)}
คือสถิติที่วัดความสอดคล้องของข้อมูลต่อเนื่องระหว่างผู้วัดตั้งแต่ 2 รายขึ้นไป
โดยพิจารณาสัดส่วนของความแปรปรวนที่เกิดจากความแตกต่างระหว่างหน่วยที่ถูกวัด (subject) เทียบกับความแปรปรวนทั้งหมด ดังสมการที่~\eqref{eq:icc}
\begin{equation}
\mathrm{ICC} = \frac{\sigma^2_{s}}{\sigma^2_{s} + \sigma^2_{r} + \sigma^2_{e}}
\label{eq:icc}
\end{equation}
โดยที่ $\sigma^2_{s}$, $\sigma^2_{r}$ และ $\sigma^2_{e}$ คือความแปรปรวนระหว่างหน่วยที่ถูกวัด ระหว่างผู้วัด และของความคลาดเคลื่อน ตามลำดับ
ICC มีหลายชนิดซึ่งให้ค่าแตกต่างกันจากข้อมูลชุดเดียวกัน จึงต้องเลือกชนิดให้เหมาะกับวัตถุประสงค์ของการวัดและระบุชนิดที่ใช้ทุกครั้ง
และแปลผลจากช่วงความเชื่อมั่นร้อยละ 95 \citep{Koo2016} ดังตารางที่~\ref{tab:agreement-interp}
ในงานวิจัยนี้ หน่วยที่ถูกวัดคือผู้ป่วย ผู้วัดคือระบบและพยาบาล และค่าที่วัดคือคะแนนรวมของแบบประเมิน
จึงใช้ ICC แบบสองทาง ชนิดความสอดคล้องสัมบูรณ์ (two-way, absolute agreement) แบบค่าเดี่ยว (single measure)
เนื่องจากชนิดความสอดคล้องสัมบูรณ์นำความแตกต่างอย่างเป็นระบบระหว่างผู้วัดมาคิดด้วย จึงตรวจสอบได้ว่าระบบให้คะแนนเท่ากับพยาบาล
ไม่ใช่เพียงให้คะแนนเปลี่ยนแปลงไปในทิศทางเดียวกัน ส่วนแบบค่าเดี่ยวสอดคล้องกับการใช้งานจริงที่คะแนนมาจากระบบเพียงแหล่งเดียว

\begin{table}[htbp]
\centering
\setlength{\tabcolsep}{6pt}
\begin{tabular}{cc@{\hspace{2.5em}}cc}
\hline
\textbf{ค่า $\kappa_w$} & \textbf{ระดับความสอดคล้อง} & \textbf{ค่า ICC} & \textbf{ระดับความสอดคล้อง} \\
\hline
น้อยกว่า 0.00 & ต่ำกว่าการสุ่ม & น้อยกว่า 0.50 & ต่ำ \\
0.00--0.20 & น้อยมาก & 0.50--0.75 & ปานกลาง \\
0.21--0.40 & พอใช้ & 0.75--0.90 & ดี \\
0.41--0.60 & ปานกลาง & มากกว่า 0.90 & ดีมาก \\
0.61--0.80 & มาก & & \\
0.81--1.00 & เกือบสมบูรณ์ & & \\
\hline
\end{tabular}
\caption{เกณฑ์การแปลผลค่า $\kappa_w$ \citep{Landis1977} และค่า ICC \citep{Koo2016}}
\label{tab:agreement-interp}
\end{table}

\subsection{งานวิจัยที่เกี่ยวข้อง}

\subsubsection{การเตรียมชุดข้อมูลบทสนทนาสำหรับการประเมิน PROMs}
\label{sec:review-dataset}

% การพัฒนาและประเมินส่วนตีความคำตอบต้องอาศัยชุดข้อมูลคำตอบของผู้ป่วยที่มีป้ายกำกับ (label) ระดับคะแนน แต่ชุดข้อมูลลักษณะนี้มักไม่มีอยู่สำหรับแบบประเมินเฉพาะทาง อีกทั้งการเก็บข้อมูลจากผู้ป่วยจริงมีต้นทุนสูงและมีข้อจำกัดด้านความเป็นส่วนตัว งานวิจัยที่เกี่ยวข้องจึงใช้แบบจำลองภาษาขนาดใหญ่สร้างข้อมูลสังเคราะห์ (synthetic data) ทดแทน

ในพัฒนาระบบ Chat-ePRO \citep{Chen2024} ทางผู้วิจัยประสบปัญหาการขาดชุดข้อมูลในการ จึงใช้ ChatGPT สวมบทบาทแพทย์เสมือนและผู้ป่วยเสมือนสนทนากันตามข้อคำถามในแบบประเมิน ข้อคำถามแต่ละข้อนิยามเป็นเนื้อหาคำถามและชุดตัวเลือก โดยเพิ่มตัวเลือก ``ไม่แน่ใจ'' และ ``ข้าม'' เพื่อรองรับกรณีพิเศษ ผู้วิจัยสร้างบทสนทนาจากแบบประเมิน 3 ฉบับรวม 4,144 ชุด แล้วให้ผู้ตรวจสอบ 2 คนคัดกรองตามเกณฑ์ 5 ข้อ เช่น บทสนทนาต้องไม่ออกนอกเรื่อง เนื้อหาต้องไม่ขัดแย้งกันเอง และต้องไม่มีข้อมูลส่วนบุคคล คงเหลือ 3,965 ชุด และกลายเป็นชุดข้อมลูสังเคราะห์ที่ใช้สำหรับฝึกแบบจำลอง
% Tang และคณะ~\citep{Tang2023} เสนอแนวทางการใช้แบบจำลองภาษาขนาดใหญ่สร้างข้อมูลสังเคราะห์พร้อมป้ายกำกับ แล้วนำไปปรับจูนแบบจำลองขนาดเล็กที่ทำงานภายในองค์กร เพื่อหลีกเลี่ยงการส่งข้อมูลผู้ป่วยไปยังบริการภายนอกและลดภาระการเก็บและติดป้ายกำกับข้อมูล ผู้วิจัยออกแบบคำสั่งสำหรับสร้างข้อมูลแบบวนซ้ำ โดยให้ ChatGPT เสนอแม่แบบคำสั่ง 5 แบบ ทดลองสร้างข้อมูลแบบละ 10 ตัวอย่าง แล้วเลือกแม่แบบที่ให้ข้อมูลคุณภาพดีที่สุดเป็นฐานของรอบถัดไป ทำซ้ำทั้งหมด 3 รอบ ในการสร้างข้อมูลแต่ละรอบจะใส่ตัวอย่างตั้งต้น (seed) ที่สุ่มจากข้อมูลจริงลงในคำสั่งเพื่อให้ข้อมูลมีความหลากหลาย และคัดตัวอย่างที่มีคุณภาพต่ำหรือซ้ำกันออกภายหลัง เมื่อทดสอบกับงานรู้จำชื่อเฉพาะทางชีวการแพทย์ (named entity recognition) และงานสกัดความสัมพันธ์ (relation extraction) แบบจำลองที่ปรับจูนด้วยข้อมูลสังเคราะห์มีค่า F1 เฉลี่ยเพิ่มขึ้นจากร้อยละ 23.37 เป็นร้อยละ 63.99 และจากร้อยละ 75.86 เป็นร้อยละ 83.59 ตามลำดับ เมื่อเทียบกับการใช้ ChatGPT แบบ zero-shot ผลการทดลองยังพบว่าการไม่ใส่ตัวอย่างตั้งต้นทำให้ ChatGPT สร้างข้อมูลซ้ำกันและประสิทธิภาพของแบบจำลองลดลงอย่างมาก และการเพิ่มจำนวนข้อมูลสังเคราะห์ช่วยเพิ่มประสิทธิภาพได้ถึงระดับหนึ่งเท่านั้น อย่างไรก็ตาม ในงานรู้จำชื่อเฉพาะ แบบจำลองที่ปรับจูนด้วยข้อมูลสังเคราะห์ยังมีประสิทธิภาพต่ำกว่าแบบจำลองที่ปรับจูนด้วยข้อมูลจริง ซึ่งผู้วิจัยอธิบายว่าเกิดจากการกระจายตัวของข้อมูลสังเคราะห์ที่แตกต่างจากข้อมูลจริง

วิธีการที่สามารถนำมาปรับใช้ในการจัดทำชุดข้อมูลคำตอบภาษาไทยสำหรับแบบประเมิน DASH และ TESS จึงมีดังนี้
\begin{enumerate}
    \item ใช้แบบจำลองภาษาขนาดใหญ่สวมบทบาทผู้ป่วยตอบข้อคำถามทีละข้อ โดยให้ตอบด้วยการบรรยายอาการหรือวิธีทำกิจกรรมแทนการระบุตัวเลือกโดยตรง และกำหนดระดับคะแนนเป้าหมายของคำตอบไว้ล่วงหน้าเพื่อใช้เป็น label
    \item สร้างคำตอบในกรณีพิเศษ เช่น คำตอบที่กำกวมและกรณีที่ผู้ป่วยไม่ได้ทำกิจกรรมนั้น และควบคุมสัดส่วนของคำตอบแต่ละระดับให้สมดุล
    % \item ใช้ตัวอย่างคำตอบตั้งต้นจำนวนหนึ่งที่ผ่านการตรวจสอบโดยบุคลากรทางการแพทย์เป็นตัวอย่างตั้งต้นในคำสั่ง เพื่อให้ข้อมูลมีความหลากหลายและใกล้เคียงภาษาพูดของผู้ป่วย ตามข้อค้นพบของ Tang และคณะ
    \item คัดกรองข้อมูลที่สร้างขึ้นตามเกณฑ์ที่กำหนด และให้ผู้ตรวจสอบตรวจทาน label
    \item ใช้ข้อมูลสังเคราะห์สำหรับการพัฒนาและเปรียบเทียบส่วนตีความคำตอบเท่านั้น ส่วนการประเมินขั้นสุดท้ายดำเนินการกับผู้ป่วยจริง เนื่องจากข้อมูลสังเคราะห์อาจมีการกระจายตัวแตกต่างจากข้อมูลจริง
\end{enumerate}

\subsubsection{การจัดการบทสนทนา (dialogue manager) สำหรับการเก็บข้อมูล PROMs}
\label{sec:review-dialogue}

การใช้ GPT-3 ควบคุมตัวแทนสนทนาเก็บข้อมูลตามช่องข้อมูล (slot) ที่ระบุในคำสั่ง
เก็บข้อมูลได้ร้อยละ 79 โดยช่องข้อมูลท้ายคำสั่งถูกข้ามบ่อยกว่าช่องข้อมูลต้นคำสั่ง และแบบจำลองไม่ได้ถามตามมาตราส่วนที่กำหนดเสมอ ~\citep{Wei2024}
และ ChatGPT สามารถติดตามสถานะของช่องข้อมูลได้แม่นยำต่ำกว่าแบบจำลองเฉพาะงาน (joint goal accuracy 0.27 เทียบกับ 0.60) แต่เมื่อได้รับสถานะที่ถูกต้อง อัตราความสำเร็จของบทสนทนาเพิ่มจาก 0.44 เป็น 0.68 ~\citep{Hudecek2023}
ส่วน CLARITY~\citep{Shaposhnikov2025} เป็นระบบคัดกรองและส่งต่อผู้ป่วยที่ใช้เครื่องสถานะจำกัด (finite-state machine) ควบคุมบทสนทนา
และใช้แบบจำลองภาษาขนาดใหญ่สร้างข้อความภายในแต่ละ state ระบบมีการใช้งานจริงกว่า 55,000 บทสนทนา และแนะนำแพทย์เฉพาะทางได้ถูกต้องในอันดับแรกร้อยละ 77

งานวิจัยนี้จึงออกแบบให้ข้อคำถามแต่ละข้อของแบบประเมินเป็นช่องข้อมูลหนึ่งช่อง และใช้ finite-state machine ในการติดตามสถานะและควบคุมลำดับการถามแทนแบบจำลอง
รวมถึงใช้เสียงที่สร้างไว้ล่วงหน้าแทนข้อความที่แบบจำลองสร้างขึ้น เพื่อไม่ให้ข้อคำถามถูกข้ามและคงถ้อยคำตามต้นฉบับ \citep{Coons2009}

\subsubsection{แบบจำลองภาษาขนาดใหญ่สำหรับการตีความคำตอบของผู้ป่วย}
\label{sec:review-interpret}

ระบบ Chat-ePRO~\citep{Chen2024} ใช้แบบจำลอง BERT ที่ปรับจูนด้วยข้อมูลสังเคราะห์ที่ได้จากการให้ ChatGPT สวมบทบาทเป็นแพทย์และคนไข้ 
สกัดตัวเลือกจากบทสนทนา ได้ค่า Macro-F1 0.90-0.95 สูงกว่า ChatGPT แบบ few-shot ที่ได้ 0.75-0.81
HopeBot~\citep{Guo2025} ใช้ GPT-4o จัดคำตอบเข้ากับระดับของแบบประเมิน PHQ-9 และคำนวณคะแนนด้วยโปรแกรมที่แยกจากแบบจำลอง
เมื่อคำตอบกำกวมจะถามเพื่อขอความชัดเจน คะแนนจากระบบสอดคล้องกับการตอบแบบประเมินด้วยตนเองที่ค่า ICC 0.92
ส่วนการเก็บแบบสำรวจทางโทรศัพท์~\citep{Kaiyrbekov2025} ใช้ GPT-4o อนุมานคำตอบจากบทถอดความและบังคับรูปแบบของผลลัพธ์
ได้ความถูกต้องร้อยละ 98 แม้บทถอดความมีอัตราความผิดพลาดระดับคำร้อยละ 7.7
อย่างไรก็ตาม ทั้งสามงานใช้แบบจำลองที่ให้บริการผ่านระบบคลาวด์ ข้อมูลของผู้ป่วยจึงต้องถูกส่งออกนอกโรงพยาบาล
ซึ่งไม่เหมาะกับข้อมูลสุขภาพตามพระราชบัญญัติคุ้มครองข้อมูลส่วนบุคคล พ.ศ. 2562 และผู้พัฒนา Chat-ePRO ระบุว่าเป็นอุปสรรคต่อการขยายผล~\citep{Chen2024}

งานวิจัยนี้จึงเลือกใช้แบบจำลองแบบ open-weight ที่เหมาะสมกับการทำงานบนเซิร์ฟเวอร์ภายใน โดยให้แบบจำลองเป็นตัวตีความคำตอบและเลือกตัวเลือกเท่านั้น 
ส่วนการคำนวณคะแนนของแบบประเมินและการจัดการบทสนทนาทำโดยส่วนอื่นของระบบ
และกำหนดให้ผลลัพธ์ระบุทั้งระดับคะแนน ค่าความมั่นใจ และเจตนาของผู้ป่วย เพื่อใช้ตัดสินว่าจะถามผู้ป่วยอีกครั้งหรือไม่
ในกรณีที่คำตอบไม่ตรงกับตัวเลือกใด หรือคำตอบไม่ชัดเจน

\subsubsection{สรุปเปรียบเทียบงานวิจัยที่เกี่ยวข้อง}
\label{sec:review-summary}

ตารางที่~\ref{tab:related-compare} สรุปเปรียบเทียบงานวิจัยที่ใช้แบบจำลองภาษาขนาดใหญ่ในการเก็บข้อมูลที่ผู้ใช้รายงานด้วยตนเองกับงานวิจัยนี้
โดยนอกจากงานที่กล่าวถึงข้างต้นแล้ว ยังรวม Talk2Care~\citep{Yang2024} ซึ่งเป็นผู้ช่วยเสียงที่ใช้ GPT-3.5 ผ่าน Amazon Alexa
เก็บข้อมูลสุขภาพจากผู้สูงอายุในรูปแบบบทสนทนาเปิดเพื่อสรุปให้บุคลากรทางการแพทย์

\begin{table}[htbp]
\centering
\small
\renewcommand{\arraystretch}{1.25}
\setlength{\tabcolsep}{3pt}
\begin{tabular}{|>{\raggedright\arraybackslash}p{2.2cm}
                |>{\raggedright\arraybackslash}p{1.7cm}
                |>{\raggedright\arraybackslash}p{2.15cm}
                |>{\raggedright\arraybackslash}p{2.2cm}
                |>{\raggedright\arraybackslash}p{2.4cm}
                |>{\raggedright\arraybackslash}p{2.5cm}|}
\hline
\textbf{งานวิจัย} & \textbf{ช่องทาง / ภาษา} & \textbf{แบบสอบถาม} & \textbf{การถามคำถาม} & \textbf{แบบจำลอง (ที่ตั้ง)} & \textbf{การประเมินผล} \\
\hline
Wei และคณะ \citep{Wei2024} & ข้อความ / อังกฤษ & ช่องข้อมูลที่ผู้วิจัยกำหนด & แบบจำลองสร้างคำถามเอง & GPT-3 (คลาวด์) & ครอบคลุมช่องข้อมูลร้อยละ 79 ($n=48$) \\
\hline
Talk2Care \citep{Yang2024} & เสียง / อังกฤษ & ข้อมูลสุขภาพแบบเปิด ไม่มีการให้คะแนน & แบบจำลองสร้างคำถามเอง & GPT-3.5 และ Amazon Alexa (คลาวด์) & ความสามารถในการใช้งาน (SUS 75.5) \\
\hline
Chat-ePRO \citep{Chen2024} & ข้อความ / จีน & Fact-B, BIBCQ และ TMI & แบบจำลองสร้างคำถามจากแบบฟอร์ม & ChatGPT (คลาวด์) และ BERT สกัดตัวเลือก & Macro-F1 0.95 บนข้อมูลสังเคราะห์ และนำร่องกับผู้ป่วย 10 คน \\
\hline
Kaiyrbekov และคณะ \citep{Kaiyrbekov2025} & เสียง (โทรศัพท์) / อังกฤษ & 33 ข้อจาก COVID Impact Survey & ตัวแทนสนทนาของ Bland AI & Bland AI และ GPT-4o (คลาวด์) & accuracy ร้อยละ 98 เทียบกับคำตอบของบุคคลสมมติ (8 คน 40 ครั้ง) \\
\hline
HopeBot \citep{Guo2025} & ข้อความหรือเสียง / อังกฤษ จีนกลาง & PHQ-9 & แบบจำลองถามด้วยภาษาทั่วไป & GPT-4o และ OpenAI TTS (คลาวด์) & ICC 0.92 เทียบกับการตอบแบบประเมินด้วยตนเอง ($n=132$) \\
\hline
\textbf{งานวิจัยนี้} & เสียง / ไทย & DASH และ TESS & เสียงตามถ้อยคำต้นฉบับที่สร้างไว้ล่วงหน้า ควบคุมลำดับด้วย finite-state machine & Typhoon Whisper Turbo และ Qwen3.5 (เซิร์ฟเวอร์ภายใน) & weighted kappa และ ICC เทียบกับการสัมภาษณ์โดยพยาบาล \\
\hline
\end{tabular}
\caption{เปรียบเทียบงานวิจัยที่ใช้แบบจำลองภาษาขนาดใหญ่ในการเก็บข้อมูลที่ผู้ใช้รายงานด้วยตนเองกับงานวิจัยนี้}
\label{tab:related-compare}
\end{table}

จากตารางที่~\ref{tab:related-compare} งานวิจัยที่ผ่านมามีข้อจำกัดร่วมกัน 3 ประการ
ประการแรก ทุกงานใช้แบบจำลองที่ให้บริการผ่านระบบคลาวด์ ซึ่งต้องส่งข้อมูลของผู้ใช้ออกไปยังผู้ให้บริการภายนอก
ประการที่สอง ทุกงานให้แบบจำลองภาษาหรือตัวแทนสนทนาสร้างถ้อยคำคำถามเอง จึงไม่สามารถรับประกันได้ว่าถ้อยคำตรงตามแบบประเมินต้นฉบับที่ผ่านการตรวจสอบความตรงแล้ว \citep{Coons2009}
ประการที่สาม ไม่มีงานใดรองรับภาษาไทย และงานที่ประเมินความสอดคล้องของคะแนนเทียบกับข้อมูลสังเคราะห์ บุคคลสมมติ หรือการตอบแบบประเมินด้วยตนเอง
ไม่ได้เทียบกับการสัมภาษณ์โดยบุคลากรทางการแพทย์ซึ่งเป็นขั้นตอนที่ใช้อยู่ในโรงพยาบาล
งานวิจัยนี้จึงพัฒนาระบบเก็บข้อมูลแบบประเมิน DASH และ TESS ด้วยเสียงภาษาไทยที่ทำงานบนเซิร์ฟเวอร์ภายในทั้งหมด
คงถ้อยคำคำถามตามต้นฉบับ และประเมินความสอดคล้องของคะแนนเทียบกับการสัมภาษณ์โดยพยาบาล
